\VerdierPhysicalPageAnchor{0025}%
where $\mathrm{RH}(X,{\ })$ denotes the total right derived functor of the
``global sections'' functor on $X$ (and likewise for $\mathrm{RH}(Y,{\ })$ and
$\mathrm{RH}(X\times Y,{\ })$), and where
$\overset{\mathrm L}{\otimes}_A$ denotes the total left derived functor of the
tensor-product functor. (We assume that $X$ and $Y$ are spaces of finite
cohomological dimension, for example finite cellular complexes.) The two
sides of $\text{(Kü)}_{\mathrm{Der}}$ are complexes of $A$-modules
determined up to quasi-isomorphism, whose cohomology modules are respectively
the abutments of the first and second spectral sequences of
$\text{(Kü)}_{\mathrm{coh}}^{4}$. The formula
$\text{(Kü)}_{\mathrm{Der}}$ is, of course, still valid when $F$ and
$G$ are complexes of sheaves on $X$ and $Y$ (suitably bounded). It is
perfectly manageable in practice, and also makes it possible to formulate
commutativity and associativity properties when there are several factors.

\VerdierSemanticAnchor{intro-5-base-change}
Extension of scalars leads to an analogous formula in derived categories.
When $A\to B$ is an $A$-algebra and $F$ is a sheaf of $A$-modules on $X$, one
has the relation
\[
\mathrm{RH}(X,F)\overset{\mathrm L}{\otimes}_A B
\xrightarrow{\sim}
\mathrm{RH}\bigl(X,F\overset{\mathrm L}{\otimes}_A B\bigr).
\]

\VerdierSemanticAnchor{intro-5-naive-formulas}
Figuratively, one may say that, in general,
\emph{naive formulas that are traditionally false become true when one works
in derived categories}.

\VerdierSemanticAnchor{intro-5-l-adic}
The importance of the coefficient rings $\mathbb Z/\ell^r\mathbb Z$ stems in
particular from the étale cohomology of schemes. It is known [6] that in this
theory one must essentially work with torsion sheaves. One obtains
$\ell$-adic cohomology, \emph{i.e.} cohomology with values in the ring
$\widehat{\mathbb Z}_{\ell}$ of $\ell$-adic integers, by taking the projective
limit of cohomology with values in the rings $\mathbb Z/\ell^r\mathbb Z$.
The formula $\text{(Kü)}_{\mathrm{Der}}$ is still valid in this
setting (when, for example, $X$ and $Y$ are proper schemes over an
algebraically closed field and $A=\mathbb Z/\ell^r\mathbb Z$), and passes
immediately to the limit to give an analogous formula for
$\widehat{\mathbb Z}_{\ell}$-sheaves. But since
$\widehat{\mathbb Z}_{\ell}$ is a principal ideal domain, the formula
$\text{(Kü)}_{\mathrm{Der}}$ gives, upon passing to cohomology, a
formula of type $\text{(Kü)}_{\mathrm{coh}}^{3}$ and, upon tensoring
with $\mathbb Q_{\ell}$, formula
$\text{(Kü)}_{\mathrm{coh}}^{1}$ for
$\mathbb Q_{\ell}$-sheaves.

\VerdierSemanticAnchor{intro-6}
\noindent\textbf{6.}\quad
It is the systematic development of the considerations sketched above that
we present in this work. Here is its plan.

\VerdierSemanticAnchor{intro-6-plan-chapter-one}
In the first chapter we study the categories of complexes of additive
categories. These categories are graded. More precisely, they are equipped
with translation functors (translation of degrees in a
